\section{The coefficient operations and their prime spectra}
\label{sec:rebuilt-spectra}
\begin{definition}
For a nonzero commutative unital ring $R$, let $G(R)=R\sqcup\{\tau\}$. Supported elements have their original ring sum and product; $x+\tau=x$ and $x\tau=\tau$. Write $e=0_R$ for the supported zero. Thus $e\ne\tau$, $e^2=e$, and $er=e$ for every supported $r$.
\end{definition}
The injection $\tau\mapsto(0,0)$, $r\mapsto(r,1)$ into $R\times\B$ proves all semiring laws, including distributivity. Here addition in $\B$ is join and multiplication is meet. The two projections are
\begin{equation}\label{eq:base-two-maps}
 p:G(R)\to R,\quad p(\tau)=0,\ p(r)=r;
 \qquad \chi:G(R)\to\B,\quad\chi(\tau)=0,\ \chi(r)=1.
\end{equation}
They jointly retain the entire original element. Every map to a ring kills $e$ since $e+e=e$, and then factors uniquely through $p$. The supported inclusion of $R$ does not preserve the specified additive zero and is not a semiring section of $p$.

\begin{theorem}[The scalar spectrum and its stalks]\label{thm:rebuilt-scalar}
The ideals of $G(R)$ are $\{\tau\}$ and $J\cup\{\tau\}$ for ring ideals $J$. Its primes are
\[
 P_\tau=\{\tau\},\qquad P_{\mathfrak p}=\mathfrak p\cup\{\tau\}\quad(\mathfrak p\in\Spec R).
\]
The map induced by $p$ identifies $\Spec R$ with $V(e)$, while $D(e)=\{P_\tau\}$ is open and dense. The localization stalks are
\[
 G(R)_{P_\tau}=\B,\qquad G(R)_{P_{\mathfrak p}}=G(R_{\mathfrak p}).
\]
For $R=\Z$, all strict maximal prime chains are
\begin{equation}\label{eq:scalar-chain}
 \{\tau\}\subsetneq\{\tau,e\}\subsetneq p\Z\cup\{\tau\}.
\end{equation}
In particular its dimension is two.
\end{theorem}
\begin{proof}
Every ideal contains $\tau$. If an ideal contains a supported element, multiplication by $e$ places $e$ in the ideal. Its supported intersection is a ring ideal because multiplication by $-1$ supplies additive inverses. This proves the ideal classification and its converse. A product equals $\tau$ only if a factor does, proving primality of $P_\tau$. Testing supported products proves primality of $J\cup\{\tau\}$ exactly when $J$ is prime.

For supported $a$, $D(a)=\{P_\tau\}\cup D_R(a)$, whereas $D(\tau)$ is empty. These formulas prove the topology assertions, and $P_\tau$ is contained in every prime. At $P_\tau$, every supported element is inverted, including $e$. Its idempotence forces $e=1$, and $er=e$ then forces every supported $r$ to equal $1$. The character $\chi$ factors through the localization and distinguishes $1$ from $\tau$, proving that the result is exactly $\B$. At $P_{\mathfrak p}$ the denominators are $R\setminus\mathfrak p$. The fraction map $r/s\mapsto r/s$, $\tau/s\mapsto\tau$ is an isomorphism: equality of supported fractions is precisely the usual localization relation; an absent fraction cannot equal a supported fraction since denominators and the extra multiplier remain supported. The list of ring primes of $\Z$ now proves \eqref{eq:scalar-chain} and excludes longer chains.
\end{proof}

\begin{proposition}[The direction of the absolute-point maps]\label{prop:base-arrows}
The maps in \eqref{eq:base-two-maps} induce
\[
 \Spec R\xrightarrow{\text{closed}}\Spec G(R)
 \xleftarrow{\text{open}}\Spec\B.
\]
There is no unital semiring map $\B\to G(R)$. The pointed multiplicative monoid $\{0,1\}$ does map to the multiplicative monoid of $G(R)$ by $0\mapsto\tau$, $1\mapsto1$.
\end{proposition}
\begin{proof}
The first assertion is the localization and prime contraction already proved. A map from $\B$ would send $1+1=1$ to $1_R+1_R=1_R$, which contradicts $1_R\ne0_R$. The two-element pointed monoid has no such additive relation, and its product table verifies the last map. Thus the two arrows and the pointed-monoid map retain their exact categories.
\end{proof}

\subsection{Every multiplicative prime, before imposing addition}
Let $M$ be the pointed multiplicative monoid underlying $G(\Z)$ and let $\mathcal P$ be the rational primes. A monoid ideal is absorbing under multiplication; no additive closure is required.
\begin{theorem}[The complete monoid spectrum]\label{thm:rebuilt-monoid}
For $\Sigma\subseteq\mathcal P$, put
\[
 J_\Sigma=\{\tau,e\}\cup\{n\in\Z\setminus\{0\}:p\mid n\text{ for some }p\in\Sigma\}.
\]
Then
\[
 \SpecMon M=\{\{\tau\}\}\sqcup\{J_\Sigma:\Sigma\subseteq\mathcal P\},
 \qquad J_\Sigma\subseteq J_{\Sigma'}\iff\Sigma\subseteq\Sigma'.
\]
This spectrum has infinite dimension. The map
\begin{equation}\label{eq:semiring-monoid-inclusion}
 j:\Spec G(\Z)\longrightarrow\SpecMon M
\end{equation}
that keeps each prime ideal as a subset is a continuous topological embedding. Its image is exactly $\{\tau\}$, $J_\varnothing$, and the $J_{\{p\}}$.
\end{theorem}
\begin{proof}
The displayed sets are proper absorbing ideals. Products in $J_\Sigma$ either involve $\tau$, vanish in the integral domain $\Z$, or have a prime factor in $\Sigma$, so they satisfy primality. If a prime monoid ideal contains a supported element it contains $e$. It contains neither unit $1,-1$. Every nonzero element in it has an integer prime factor in it, by repeated primality and integer factorization. Conversely it contains every multiple of every such prime factor. This proves exhaustion, including $J_\varnothing=\{\tau,e\}$. Membership of a rational prime $p$ tests whether $p\in\Sigma$, proving the inclusion statement. Arbitrarily long strictly nested finite subsets give arbitrarily long prime chains.

The scalar theorem identifies the image of $j$. For every $a\in M$, the inverse image under $j$ of the monoid basic open $D_M(a)$ is the semiring basic open $D_{G(\Z)}(a)$. These basics also generate the topology on the source, proving that the injection is an embedding. Finally if $\Sigma$ contains distinct primes $p,q$, B\'ezout integers $u,v$ give $up+vq=1$. Both summands belong to $J_\Sigma$, so additive closure would force $1$ into it. Such a monoid prime therefore cannot be a semiring prime. This also proves exactly which additive relation removes it.
\end{proof}

\begin{proposition}[Reconstruction retains the missing relations]\label{prop:preaddition-rebuild}
Let $\N_0[M]$ be the monoid semiring with the basis symbol $[\tau]$ identified with its additive zero. Then the quotient by
\[
 [a]+[b]\sim[a+b]\qquad(a,b\in\Z)
\]
is canonically $G(\Z)$.
\end{proposition}
\begin{proof}
Evaluation $[a]\mapsto a$ respects all displayed relations. Every nonempty supported sum reduces to the single basis symbol of its integer sum, while the empty sum remains the symbol $\tau$. Conversely $a\mapsto[a]$, $\tau\mapsto0$ preserves multiplication, every supported addition, and addition with the specified zero. These maps are inverse. The contraction of $[\tau]$ is essential: in the uncontracted free semiring, the relations $[x]+[\tau]=[x]$ alone do not identify $[\tau]$ with the empty sum. Indeed they hold in the semiring obtained by adjoining another external zero below the already specified zero of $G(\Z)$, where the two remain distinct. This gives an explicit separating model for the missing relation.
\end{proof}

\subsection{The complete lattice of zero-amplitude support}
Let $L$ be a nontrivial bounded distributive lattice and put
\[
 S_L=G_L(R)=\{(0,\lambda):\lambda\in L\}\cup\{(r,1_L):r\in R\}
 \subset R\times L.
\]
Operations are $(r,\lambda)+(s,\mu)=(r+s,\lambda\vee\mu)$ and
$(r,\lambda)(s,\mu)=(rs,\lambda\wedge\mu)$. Set $z_\lambda=(0,\lambda)$, $\tau_L=z_0$, $e_L=z_1$, and $Z_L=\{z_\lambda\}$.
\begin{theorem}[All support and arithmetic primes]\label{thm:rebuilt-full-support}
The primes of $S_L$ are exactly
\[
 P_{\mathfrak a}=\{z_\lambda:\lambda\in\mathfrak a\},\quad\mathfrak a\in\SpecLat L;
 \qquad Q_{\mathfrak p}=Z_L\cup\{(r,1):r\in\mathfrak p\},\quad\mathfrak p\in\Spec R.
\]
Within each family, inclusions retain the original prime order; every $P_{\mathfrak a}$ is strictly contained in every $Q_{\mathfrak p}$. Topologically
\begin{equation}\label{eq:full-spectral-base}
 \Spec S_L=\SpecLat L\blacktriangleright\Spec R,
 \qquad D(e_L)=\SpecLat L,\quad V(e_L)=\Spec R,
\end{equation}
where opens are opens of the first stratum and sets consisting of that entire stratum together with an open of the second. If both dimensions are finite, their dimensions add with one further step between the strata.
\end{theorem}
\begin{proof}
An ideal containing a supported element contains $e_L$ by multiplication by it and therefore all of $Z_L$, since $e_Lz_\lambda=z_\lambda$. Its supported elements form a ring ideal. Every remaining ideal is $\{z_\lambda:\lambda\in\mathfrak a\}$ for a proper lattice ideal: closure under join follows from addition, and downward closure follows from multiplication by $z_\mu$ for $\mu\le\lambda$. Conversely these are all ideals. Primality in the first family is the meet-prime condition on $\mathfrak a$; in the second it is the ring-prime condition. Products involving a zero-fibre element obey exactly these tests. The inclusions follow from the displayed sets. Finally
\[
 D(z_\lambda)=\{P_{\mathfrak a}:\lambda\notin\mathfrak a\},\qquad
 D((r,1))=\SpecLat L\cup D_R(r)
\]
prove the topology and the two strata. Every chain is a support chain followed by an arithmetic chain, and maximal finite chains concatenate, proving the dimension formula.
\end{proof}

\begin{proposition}[Two support channels and their exact comparison]\label{prop:two-support-branches}
For $L=\mathcal P(\{1,2\})$, $\{\tau_L\}$ is not prime, since $z_{\{1\}}z_{\{2\}}=\tau_L$ although neither factor is absent. There are two support primes
\[
 P_1=\{\tau_L,z_{\{2\}}\},\qquad P_2=\{\tau_L,z_{\{1\}}\}.
\]
For $R=\Z$, each specializes to the common $Q_{(0)}$ and then to every $Q_{(p)}$. Moreover
\begin{equation}\label{eq:common-amplitude-fibre-product}
 G_{\B^2}(R)\simeq G(R)\mathbin{\times_R}G(R),
\end{equation}
where the two maps to $R$ are ring reflections. The two factor-generic primes contract to $P_1,P_2$, while their arithmetic primes above the same $\mathfrak p$ both contract to $Q_{\mathfrak p}$.
\end{proposition}
\begin{proof}
The two coordinate maps $L\to\B$ have precisely the displayed prime kernels; a proper prime ideal must contain exactly one of the two atoms, proving exhaustion. For \eqref{eq:common-amplitude-fibre-product}, send $z_A$ to the pair with entry $e$ on $A$ and $\tau$ off $A$, and send $(r,\{1,2\})$ to $(r,r)$. Equal nonzero amplitudes force both coordinates to be that supported value, whereas common amplitude zero allows exactly the four masks. Hence the map is onto the fibre product and injective; the operations verify it is an isomorphism. Taking inverse images of the factor primes proves their contractions.
\end{proof}

\begin{proposition}[What a Boolean branch retains]\label{prop:branch-contraction}
A bounded lattice homomorphism $\epsilon:L\to\B$ induces $\beta_\epsilon:G_L(R)\to G_\B(R)$ by $z_\lambda\mapsto z_{\epsilon(\lambda)}$ and $(r,1)\mapsto(r,1)$. Its spectral image is
\[
 \{P_{\epsilon^{-1}(0)}\}\blacktriangleright\Spec R.
\]
\end{proposition}
\begin{proof}
The map preserves joins, meets and supported arithmetic, hence is a semiring homomorphism. The inverse image of the scalar absent prime is $\{z_\lambda:\epsilon(\lambda)=0\}$, and the inverse image of each arithmetic prime contains the full zero fibre and exactly its original amplitudes. This proves the asserted image. In particular a branch selects one support prime; the formula does not recover the other primes.
\end{proof}

These proofs reconstruct the scalar and multiplicative results of \cite{globalization,secondary} and the lattice-support results of \cite{support}. They give explicit maps whenever their spectra differ. The following section uses the multiplicative prime faces to recover the finite sets of places occurring in the arithmetic trace, retaining the support geometry of \eqref{eq:full-spectral-base} for the cohomological coefficient construction.
